% !TEX root = ../main.tex
% --- LỜI CẢM ƠN ---
\newpage
\pagestyle{frontmatter}
\chapter*{LỜI CẢM ƠN}
\thispagestyle{frontmatter}

Bốn năm đại học tưởng chừng còn dài phía trước, vậy mà nay đã đến lúc chúng em ngồi viết những dòng cuối cùng cho khóa luận tốt nghiệp của mình. Với chúng em, đây không chỉ đơn thuần là một bản báo cáo khép lại quãng đời sinh viên, mà còn là kết tinh của biết bao đêm thức trắng bên màn hình, những lần loay hoay gỡ từng dòng mã nguồn đến rạng sáng, và cả niềm vui vỡ òa khi sản phẩm dần thành hình. Nhìn lại cả chặng đường ấy, chúng em hiểu rằng nỗ lực của riêng mình là chưa đủ; để đi được đến ngày hôm nay, chúng em đã may mắn nhận được sự dìu dắt, sẻ chia và tin tưởng từ rất nhiều người.

Trước hết, với tất cả sự kính trọng và lòng biết ơn sâu sắc nhất, chúng em xin gửi lời cảm ơn chân thành đến cô \textbf{ThS. Nguyễn Thị Hoàng Khánh} --- người đã trực tiếp hướng dẫn chúng em trong suốt quá trình thực hiện đề tài. Cô không chỉ định hình lại cho chúng em những ý tưởng còn non nớt buổi ban đầu, mà còn kiên nhẫn lắng nghe, tỉ mỉ góp ý và thẳng thắn chỉ ra những thiếu sót mà tự chúng em khó lòng nhận ra. Có những lúc cả nhóm rơi vào bế tắc, hoang mang không biết nên đi tiếp theo hướng nào, chính sự tận tâm và những lời động viên đúng lúc của Cô đã tiếp thêm cho chúng em niềm tin để bước tiếp. Những kinh nghiệm quý báu mà Cô trao gửi sẽ là hành trang theo chúng em đi rất xa trên chặng đường phía trước.

Chúng em cũng xin bày tỏ lòng biết ơn đến Quý Thầy Cô Khoa Công nghệ Thông tin nói riêng và toàn thể Thầy Cô Trường Đại học Công nghiệp Thành phố Hồ Chí Minh nói chung. Từng bài giảng, từng buổi thực hành, và cả những lời nhắc nhở nghiêm khắc của Thầy Cô trong suốt bốn năm qua đã trang bị cho chúng em không chỉ kiến thức chuyên môn vững vàng, mà còn cả tư duy giải quyết vấn đề và tinh thần tự học --- những thứ mà chúng em tin rằng sẽ còn nguyên giá trị rất lâu sau khi rời khỏi giảng đường. Chúng em cũng xin cảm ơn Ban Giám hiệu Nhà trường đã luôn tạo điều kiện thuận lợi về cơ sở vật chất cùng môi trường học tập và nghiên cứu, để chúng em được yên tâm theo đuổi đam mê của mình.

Chúng em xin chân thành cảm ơn!

\vspace{1cm}
\begin{flushright}
    \textit{TP. Hồ Chí Minh, ngày 20 tháng 4 năm 2026}\\
    \textbf{Sinh viên thực hiện}\\
    \textbf{Lai Thanh Sĩ}\\
    \textbf{Phạm Đức Tài}
\end{flushright}

% --- NHẬN XÉT GVHD ---
\newpage
\chapter*{NHẬN XÉT CỦA GIẢNG VIÊN HƯỚNG DẪN}
\thispagestyle{frontmatter}
\dotline \dotline \dotline \dotline \dotline
\dotline \dotline \dotline \dotline \dotline
\dotline \dotline \dotline 
\vspace{0.5cm}
\begin{flushright}
    \textit{TP. Hồ Chí Minh, ngày \dots tháng \dots năm 20\dots}\\
    \textbf{GIẢNG VIÊN HƯỚNG DẪN}\\
    \textit{(Ký và ghi rõ họ tên)}\\
    \vspace{1.5cm}
\end{flushright}

% --- NHẬN XÉT GVPB 1 ---
\newpage
\chapter*{NHẬN XÉT CỦA GIẢNG VIÊN PHẢN BIỆN 1}
\thispagestyle{frontmatter}
\dotline \dotline \dotline \dotline \dotline
\dotline \dotline \dotline \dotline \dotline
\dotline \dotline \dotline 
\vspace{0.5cm}
\begin{flushright}
    \textit{TP. Hồ Chí Minh, ngày \dots tháng \dots năm 20\dots}\\
    \textbf{GIẢNG VIÊN PHẢN BIỆN 1}\\
    \textit{(Ký và ghi rõ họ tên)}\\
    \vspace{1.5cm}
\end{flushright}

% --- NHẬN XÉT GVPB 2 ---
\newpage
\chapter*{NHẬN XÉT CỦA GIẢNG VIÊN PHẢN BIỆN 2}
\thispagestyle{frontmatter}
\dotline \dotline \dotline \dotline \dotline
\dotline \dotline \dotline \dotline \dotline
\dotline \dotline \dotline 
\vspace{0.5cm}
\begin{flushright}
    \textit{TP. Hồ Chí Minh, ngày \dots tháng \dots năm 20\dots}\\
    \textbf{GIẢNG VIÊN PHẢN BIỆN 2}\\
    \textit{(Ký và ghi rõ họ tên)}\\
    \vspace{1.5cm}
\end{flushright}

% --- MỤC LỤC ---
\clearpage
\pagestyle{frontmatter}
\begingroup
\let\clearpage\relax
\addcontentsline{toc}{chapter}{MỤC LỤC}
\tableofcontents
\endgroup
\thispagestyle{frontmatter}

% --- DANH MỤC HÌNH ẢNH ---
\clearpage
\begingroup
\let\clearpage\relax
\addcontentsline{toc}{chapter}{DANH MỤC CÁC HÌNH ẢNH}
\listoffigures
\endgroup
\thispagestyle{frontmatter}

% --- DANH MỤC BẢNG BIỂU ---
\clearpage
\begingroup
\let\clearpage\relax
\addcontentsline{toc}{chapter}{DANH MỤC CÁC BẢNG BIỂU}
\listoftables
\endgroup
\thispagestyle{frontmatter}

% --- DANH MỤC TỪ VIẾT TẮT ---
\clearpage
\chapter*{DANH MỤC CÁC TỪ VIẾT TẮT}
\addcontentsline{toc}{chapter}{DANH MỤC CÁC TỪ VIẾT TẮT}
\thispagestyle{frontmatter}

% Chia tỷ lệ các cột: STT(8) + Từ viết tắt(18) + Nghĩa đầy đủ(74) = 100\tbu
\setlength{\tbu}{\dimexpr(\linewidth - 6\tabcolsep - 4\arrayrulewidth)/100\relax}
\begin{longtable}{|>{\centering\arraybackslash}p{8\tbu}|p{18\tbu}|p{74\tbu}|}
    \hline
    \textbf{STT} & \textbf{Từ viết tắt} & \textbf{Nghĩa đầy đủ} \\
    \hline
    \endfirsthead
    \hline
    \textbf{STT} & \textbf{Từ viết tắt} & \textbf{Nghĩa đầy đủ} \\
    \hline
    \endhead
    1  & AES   & Automated Essay Scoring (Chấm điểm bài viết tự động) \\
    \hline
    2  & AI    & Artificial Intelligence (Trí tuệ nhân tạo) \\
    \hline
    3  & AMQP  & Advanced Message Queuing Protocol \\
    \hline
    4  & AMQPS & Advanced Message Queuing Protocol Secure (AMQP qua TLS) \\
    \hline
    5  & API   & Application Programming Interface \\
    \hline
    6  & APS   & Automated Pronunciation Scoring (Chấm điểm phát âm tự động) \\
    \hline
    7  & BaaS  & Backend as a Service \\
    \hline
    8  & CALL  & Computer-Assisted Language Learning (Học ngôn ngữ hỗ trợ máy tính) \\
    \hline
    9  & CDN   & Content Delivery Network (Mạng phân phối nội dung) \\
    \hline
    10 & CEFR  & Common European Framework of Reference for Languages \\
    \hline
    11 & CI/CD & Continuous Integration / Continuous Deployment (Tích hợp và triển khai liên tục) \\
    \hline
    12 & DI    & Dependency Injection (Tiêm phụ thuộc) \\
    \hline
    13 & DLQ   & Dead Letter Queue (Hàng đợi thư chết) \\
    \hline
    14 & DTO   & Data Transfer Object \\
    \hline
    15 & ERD   & Entity--Relationship Diagram (Sơ đồ thực thể -- quan hệ) \\
    \hline
    16 & FSRS  & Free Spaced Repetition Scheduler \\
    \hline
    17 & GCP   & Google Cloud Platform \\
    \hline
    18 & GCR   & Google Container Registry (Kho lưu trữ container của Google) \\
    \hline
    19 & GCS   & Google Cloud Storage \\
    \hline
    20 & HMAC  & Hash-based Message Authentication Code \\
    \hline
    21 & HTTP  & HyperText Transfer Protocol \\
    \hline
    22 & HTTPS & HyperText Transfer Protocol Secure \\
    \hline
    23 & IDOR  & Insecure Direct Object Reference (Tham chiếu đối tượng trực tiếp không an toàn) \\
    \hline
    24 & IELTS & International English Language Testing System \\
    \hline
    25 & IPA   & International Phonetic Alphabet (Bảng phiên âm quốc tế) \\
    \hline
    26 & IPN   & Instant Payment Notification (Thông báo thanh toán tức thời) \\
    \hline
    27 & JSON  & JavaScript Object Notation \\
    \hline
    28 & JWT   & JSON Web Token \\
    \hline
    29 & LLM   & Large Language Model (Mô hình ngôn ngữ lớn) \\
    \hline
    30 & MAU   & Monthly Active Users (Người dùng hoạt động hàng tháng) \\
    \hline
    31 & ML    & Machine Learning (Học máy) \\
    \hline
    32 & NLP   & Natural Language Processing (Xử lý ngôn ngữ tự nhiên) \\
    \hline
    33 & OOP   & Object-Oriented Programming (Lập trình hướng đối tượng) \\
    \hline
    34 & ORM   & Object-Relational Mapping \\
    \hline
    35 & OTLP  & OpenTelemetry Protocol \\
    \hline
    36 & RDBMS & Relational Database Management System (Hệ quản trị cơ sở dữ liệu quan hệ) \\
    \hline
    37 & REST  & Representational State Transfer \\
    \hline
    38 & SDK   & Software Development Kit \\
    \hline
    39 & SQL   & Structured Query Language \\
    \hline
    40 & SRS   & Spaced Repetition System (Hệ thống lặp lại ngắt quãng) \\
    \hline
    41 & SSG   & Static Site Generation \\
    \hline
    42 & SSH   & Secure Shell (Giao thức kết nối đầu cuối an toàn) \\
    \hline
    43 & SSL   & Secure Sockets Layer \\
    \hline
    44 & SSR   & Server-Side Rendering \\
    \hline
    45 & STT   & Speech-to-Text (Nhận dạng giọng nói thành văn bản) \\
    \hline
    46 & TLS   & Transport Layer Security \\
    \hline
    47 & TTL   & Time to Live \\
    \hline
    48 & UC    & Use Case (Ca sử dụng) \\
    \hline
    49 & UI    & User Interface (Giao diện người dùng) \\
    \hline
    50 & UUID  & Universally Unique Identifier (Mã định danh duy nhất toàn cầu) \\
    \hline
    51 & UX    & User Experience (Trải nghiệm người dùng) \\
    \hline
    52 & VM    & Virtual Machine (Máy ảo) \\
    \hline
    53 & XP    & Experience Points (Điểm kinh nghiệm) \\
    \hline
    54 & ZPD   & Zone of Proximal Development (Vùng phát triển gần) \\
    \hline
\end{longtable}
